\section{Introduction}\label{sec:introduction}
\COtwo solubility in aqueous alkanolamines is a central thermodynamic quantity for carbon-capture solvents.  Monoethanolamine (\MEA) remains a reference solvent because of its industrial history and extensive experimental literature.  Absorbed \COtwo reacts with water and amine to form molecular \COtwo, protonated amine, carbamate, bicarbonate, carbonate, hydronium, and hydroxide.  The equilibrium \COtwo pressure at a given loading consequently depends on both phase equilibrium and liquid speciation.

Reactive vapor--liquid-equilibrium measurements commonly express \COtwo solubility as partial pressure versus loading.  Pressure measurements test the gas--liquid fugacity balance, while spectroscopic composition measurements test reaction equilibria, material balances, and electroneutrality.  Fitting pressure alone can conceal an implausible distribution of \MEA, \MEAH, \MEACOO, \HCOthree, and \COthree.  We therefore evaluate pressure and speciation as coupled observations of the same thermodynamic state.

\subsection{CO2 solubility in reactive MEA}
Experimental data often specify unloaded amine mass fraction and absorbed \COtwo per mole of amine.  Material and charge balances map these apparent variables to a true-species liquid composition before the equation of state evaluates fugacities and activity coefficients.  The evidence used here combines vapor--liquid-equilibrium (VLE) measurements of \COtwo partial pressure with nuclear magnetic resonance (NMR) measurements of liquid species.  VLE data define the pressure response, and NMR data constrain the underlying chemical composition.

\subsection{Thermodynamic modeling of CO2 solubility in MEA--CO2--H2O}
Thermodynamic descriptions of \COtwo-loaded amines have historically relied on apparent-species correlations or activity-coefficient models.  Kent--Eisenberg and related apparent-equilibrium approaches provide compact correlations with transferability limited to their fitted temperature, loading, and solvent-composition windows.  Electrolyte-NRTL and related excess-Gibbs-energy models represent ionic liquid nonideality for \MEA--\COtwo--\HtwoO systems\cite{Zhang2011,Morgan2017b,Mouhoubi2020,Hilliard2008}.  Their reliability depends on fitted interaction parameters and the species convention used for regression.

The \MEA literature contains pressure and speciation measurements suitable for testing a molecularly explicit description of \COtwo solubility.  VLE studies report \COtwo partial pressures across temperature, loading, and amine-composition conditions\cite{Jou1995,Hilliard2008,Mamun2005,Xu2011,Aronu2011}.  NMR studies report liquid-phase \MEA, \MEAH, \MEACOO, \HCOthree, and carbonate-containing observables at several levels of species resolution\cite{Jakobsen2005,Bottinger2008,Matin2012}.  A reactive-solubility model can thus be assessed against both gas pressure and liquid composition.

\subsection{SAFT and PC-SAFT models for CO2 solubility}
Statistical associating fluid theory (SAFT)-type equations of state provide an alternative to purely excess-Gibbs-energy descriptions because molecular size, dispersion, chain formation, and association are represented in the Helmholtz-energy model.  The perturbed-chain statistical associating fluid theory (PC-SAFT) equation of Gross and Sadowski introduced a perturbation theory for chain molecules\cite{Gross2001}, and its associating extension provides the site-based hydrogen-bonding contribution needed for water and alkanolamines\cite{Gross2002}.  PC-SAFT has been applied directly to aqueous ethanolamine systems and acid-gas solubility, including pure alkanolamine property regression and VLE prediction for aqueous amines\cite{Nasrifar2010}.  Baygi and Pahlavanzadeh applied PC-SAFT to \COtwo solubility in aqueous \MEA, using a chemical-equilibrium step to reconstruct liquid species and a neutral PC-SAFT pressure calculation as the phase-equilibrium model\cite{Baygi2015}.  Other SAFT variants have also been used for loaded \MEA, including transferable SAFT-VR models that represent chemisorption stoichiometry through effective association sites\cite{MacDowell2010}.  More recent SAFT comparisons for alkanolamines continue to show that EOS-based models can describe pure, binary, and ternary amine properties with physically interpretable parameters\cite{Pakravesh2025a}.

Most PC-SAFT treatments of loaded \MEA use an apparent chemistry representation or solve the reactive liquid composition outside the electrolyte EOS.  They then pass that composition to the phase-equilibrium calculation instead of deriving ionic activity coefficients from the same EOS.  This separation matters for \MEA because carbamate formation creates the \MEAH/\MEACOO ion pair absent from tertiary-amine MDEA chemistry.

\subsection{ePC-SAFT and advanced Born electrostatics}
Electrolyte PC-SAFT (ePC-SAFT) extends PC-SAFT to charged species by adding an electrolyte contribution to the residual Helmholtz energy.  The original aqueous-electrolyte formulation introduced Debye--Huckel electrostatic effects into the PC-SAFT framework and demonstrated correlations for aqueous salt densities and vapor pressures\cite{Cameretti2005}.  Uyan et al. applied ePC-SAFT to aqueous MDEA--\COtwo solubility with explicit ionic species and activity-coupled chemical equilibrium, treating the calculation as predictive when no binary parameters were adjusted to the ternary acid-gas data\cite{Uyan2015}.  Wangler et al. extended that amine-solubility strategy by fitting selected binary interactions to independent osmotic-coefficient data and validating the resulting model against acid-gas solubility, thereby distinguishing literature-collected pure-component parameters, independently regressed binary parameters, and final VLE validation data\cite{Wangler2018}.

Later ePC-SAFT developments addressed non-aqueous and mixed-solvent electrolytes, where ion solvation and composition-dependent permittivity become central.  The ePC-SAFT advanced formulation of Buelow, Ascani, and Held incorporated a concentration-dependent dielectric constant into both the Born contribution and Debye--Huckel term\cite{Bulow2020}; the second part extended the framework to salt solubility and ion pairing\cite{Bulow2021}.  Subsequent work used ePC-SAFT advanced to predict \COtwo solubility in aqueous and organic electrolyte solutions, including systems with carbonate chemistry\cite{Schick2023}.  More recent studies refined the modified-Born term through solvation-shell modeling and dielectric saturation (SSM+DS), which add ion-local solvation structure to the Born contribution\cite{Rueben2024,Figiel2025}.

Loaded aqueous \MEA tests this model class because \COtwo loading changes the ionic population, solvent dielectric environment, and carbamate concentration.  The system combines associating neutral molecules, molecular \COtwo, changing solvent composition, and several ionic products of acid--base chemistry.  A true-species SSM+DS ePC-SAFT model links phase equilibrium, reaction equilibrium, ionic activity coefficients, and composition-dependent solvation within one thermodynamic framework.

\begin{table*}
    \centering
    \scriptsize
    \caption{Positioning of the present \MEA \COtwo-solubility study relative to representative SAFT-family and electrolyte-model studies.}
    \label{tab:literature-model-comparison}
    \begin{tabularx}{\textwidth}{@{}>{\raggedright\arraybackslash}p{2.2cm}>{\raggedright\arraybackslash}X>{\raggedright\arraybackslash}X>{\raggedright\arraybackslash}X>{\raggedright\arraybackslash}X@{}}
        \toprule
        Study & System and model & Species and activity treatment & Data and fitted quantities & Validation reported \\
        \midrule
        Nasrifar and Tafazzol\cite{Nasrifar2010} & \MEA--\HtwoO and acid-gas PC-SAFT mixtures & Apparent neutral mixture basis without explicit Born treatment & Pure-component and VLE/property data; molecular parameters and binary interactions & Pure and mixture thermodynamic-property checks \\
        Baygi and Pahlavanzadeh\cite{Baygi2015} & \MEA--\HtwoO--\COtwo PC-SAFT with external chemistry & Six-species reconstructed liquid without ionic ePC-SAFT activity coupling or SSM+DS Born terms & \COtwo pressure data with reconstructed speciation; neutral PC-SAFT molecular and binary terms & Pressure curves and reconstructed speciation \\
        Uyan et al.\cite{Uyan2015} & MDEA--\HtwoO--\COtwo ePC-SAFT & Explicit ionic true-species chemistry with electrolyte activity terms; no \MEA carbamate ion pair & Solubility data and independent parameter sources; selected ion and binary parameters & Predictive ternary acid-gas solubility checks \\
        Wangler et al.\cite{Wangler2018} & MDEA amine ePC-SAFT systems & Explicit ionic true species with fitted electrolyte binary interactions; no \MEACOO pathway & Osmotic, solubility, and related data; selected binary and pure-ion quantities & Acid-gas solubility and related validation data \\
        Buelow et al.\cite{Bulow2020,Bulow2021} & Electrolyte and mixed-solvent advanced ePC-SAFT & Concentration-dependent permittivity with advanced Born corrections & Electrolyte solution and salt-solubility data; dielectric/Born framework terms & Electrolyte and salt-solubility validation \\
        Figiel et al.\cite{Figiel2025} & Electrolyte and organic-solvent modified-Born ePC-SAFT & Explicit ionic true species with SSM+DS modified Born treatment & Electrolyte activity and solubility data; Born, dielectric, and solvent-shell conventions & Activity, solubility, and consistency checks \\
        This work & \MEA--\HtwoO--\COtwo SSM+DS modified-Born ePC-SAFT & Nine-species explicit ionic basis with \MEAH/\MEACOO carbamate chemistry, SSM+DS Born terms, and concentration-dependent permittivity & 161 \COtwo-solubility pressure rows and 74 NMR-derived speciation rows; direct \MEAH/\MEACOO ion regression and fixed-parameter activity evaluation & Solubility-pressure/speciation residuals, train-validation split, and carbonate Born identifiability screen \\
        \bottomrule
    \end{tabularx}
\end{table*}


\subsection{Contribution of this work}
We develop a nine-species SSM+DS modified-Born ePC-SAFT model for reactive aqueous \MEA--\HtwoO--\COtwo.  The calculation couples reaction equilibria to ePC-SAFT activities, fugacities, concentration-dependent permittivity, and modified-Born solvation.  We fit the \MEAH/\MEACOO parameters to NMR speciation measurements, then evaluate the resulting model against additional NMR states and 161 \COtwo partial-pressure records.  The comparison tests both the equilibrium pressure and the liquid species distribution predicted by one thermodynamic state.
